% Author biographies and photos, kept out of the submission: IEEE asks for them only after acceptance.
% To restore, uncomment % Author biographies and photos, kept out of the submission: IEEE asks for them only after acceptance.
% To restore, uncomment % Author biographies and photos, kept out of the submission: IEEE asks for them only after acceptance.
% To restore, uncomment % Author biographies and photos, kept out of the submission: IEEE asks for them only after acceptance.
% To restore, uncomment \input{bios} in main.tex (after \bibliography).
% Author biographies (supervisor, 2026-11-24). Photos in manuscript/photos/ (git-ignored: the repo is public).
% Leu: degrees and career from his NTUST profile (et.ntust.edu.tw, fetched 2026-11-23) and the biography in his
% IEEE papers; photo from the ICOMPAC 2025 keynote page (icompac.its.ac.id), the largest found (1101x1357).
% Ardiansyah: from the owner's CV. Yunanto: from the owner (2026-11-24).
\begin{IEEEbiography}[{\includegraphics[width=1in,height=1.25in,clip,keepaspectratio]{photos/yunanto.jpg}}]{Rifky Ari Yunanto}
is pursuing the M.S. degree with the College of Electrical Engineering and Computer Science, National Taiwan University of Science and Technology (NTUST), Taipei, Taiwan, where he is a member of the Mobilizing Information Technology Laboratory (MITLab).
% TODO: Rifky's bachelor's degree (field, university, year) and research interests.
\end{IEEEbiography}

\begin{IEEEbiography}[{\includegraphics[width=1in,height=1.25in,clip,keepaspectratio]{photos/leu.jpg}}]{Jenq-Shiou Leu}
(Senior Member, IEEE) received the B.S. degree in mathematics and the M.S. degree in computer science and information engineering from National Taiwan University, Taipei, Taiwan, and the Ph.D. degree in computer science from National Tsing Hua University, Hsinchu, Taiwan.
From 1995 to 2007 he worked in industry, with Rising Star Technology, Mobitai Communications and Taiwan Mobile.
He joined the Department of Electronic and Computer Engineering, National Taiwan University of Science and Technology (NTUST), Taipei, Taiwan, in 2007, where he has been a Professor since 2014, chaired the department from 2017 to 2020, and has been Dean of the College of Electrical Engineering and Computer Science since 2024.
He leads the Mobilizing Information Technology Laboratory (MITLab).
His research interests include B5G/6G technology, AI/ML platforms, applications of large language models and biomedical information analysis.
\end{IEEEbiography}

\begin{IEEEbiography}[{\includegraphics[width=1in,height=1.25in,clip,keepaspectratio]{photos/ardiansyah.jpg}}]{Dhista Dwi Nur Ardiansyah}
received the bachelor's degree in informatics engineering education from Universitas Negeri Yogyakarta, Yogyakarta, Indonesia, in 2024.
Since 2026, Ardiansyah has been a Research Assistant with the Mobilizing Information Technology Laboratory (MITLab), National Taiwan University of Science and Technology, Taipei, Taiwan, through the Taiwan Experience Education Program (TEEP), working on small language models on edge devices, and has worked as a software engineer in Indonesia since 2021.
As an undergraduate, Ardiansyah won a gold medal at the 2021 Indonesian student digital innovation competition (LIDM) for an automatic question generator for teachers.
\end{IEEEbiography}
 in main.tex (after \bibliography).
% Author biographies (supervisor, 2026-11-24). Photos in manuscript/photos/ (git-ignored: the repo is public).
% Leu: degrees and career from his NTUST profile (et.ntust.edu.tw, fetched 2026-11-23) and the biography in his
% IEEE papers; photo from the ICOMPAC 2025 keynote page (icompac.its.ac.id), the largest found (1101x1357).
% Ardiansyah: from the owner's CV. Yunanto: from the owner (2026-11-24).
\begin{IEEEbiography}[{\includegraphics[width=1in,height=1.25in,clip,keepaspectratio]{photos/yunanto.jpg}}]{Rifky Ari Yunanto}
is pursuing the M.S. degree with the College of Electrical Engineering and Computer Science, National Taiwan University of Science and Technology (NTUST), Taipei, Taiwan, where he is a member of the Mobilizing Information Technology Laboratory (MITLab).
% TODO: Rifky's bachelor's degree (field, university, year) and research interests.
\end{IEEEbiography}

\begin{IEEEbiography}[{\includegraphics[width=1in,height=1.25in,clip,keepaspectratio]{photos/leu.jpg}}]{Jenq-Shiou Leu}
(Senior Member, IEEE) received the B.S. degree in mathematics and the M.S. degree in computer science and information engineering from National Taiwan University, Taipei, Taiwan, and the Ph.D. degree in computer science from National Tsing Hua University, Hsinchu, Taiwan.
From 1995 to 2007 he worked in industry, with Rising Star Technology, Mobitai Communications and Taiwan Mobile.
He joined the Department of Electronic and Computer Engineering, National Taiwan University of Science and Technology (NTUST), Taipei, Taiwan, in 2007, where he has been a Professor since 2014, chaired the department from 2017 to 2020, and has been Dean of the College of Electrical Engineering and Computer Science since 2024.
He leads the Mobilizing Information Technology Laboratory (MITLab).
His research interests include B5G/6G technology, AI/ML platforms, applications of large language models and biomedical information analysis.
\end{IEEEbiography}

\begin{IEEEbiography}[{\includegraphics[width=1in,height=1.25in,clip,keepaspectratio]{photos/ardiansyah.jpg}}]{Dhista Dwi Nur Ardiansyah}
received the bachelor's degree in informatics engineering education from Universitas Negeri Yogyakarta, Yogyakarta, Indonesia, in 2024.
Since 2026, Ardiansyah has been a Research Assistant with the Mobilizing Information Technology Laboratory (MITLab), National Taiwan University of Science and Technology, Taipei, Taiwan, through the Taiwan Experience Education Program (TEEP), working on small language models on edge devices, and has worked as a software engineer in Indonesia since 2021.
As an undergraduate, Ardiansyah won a gold medal at the 2021 Indonesian student digital innovation competition (LIDM) for an automatic question generator for teachers.
\end{IEEEbiography}
 in main.tex (after \bibliography).
% Author biographies (supervisor, 2026-11-24). Photos in manuscript/photos/ (git-ignored: the repo is public).
% Leu: degrees and career from his NTUST profile (et.ntust.edu.tw, fetched 2026-11-23) and the biography in his
% IEEE papers; photo from the ICOMPAC 2025 keynote page (icompac.its.ac.id), the largest found (1101x1357).
% Ardiansyah: from the owner's CV. Yunanto: from the owner (2026-11-24).
\begin{IEEEbiography}[{\includegraphics[width=1in,height=1.25in,clip,keepaspectratio]{photos/yunanto.jpg}}]{Rifky Ari Yunanto}
is pursuing the M.S. degree with the College of Electrical Engineering and Computer Science, National Taiwan University of Science and Technology (NTUST), Taipei, Taiwan, where he is a member of the Mobilizing Information Technology Laboratory (MITLab).
% TODO: Rifky's bachelor's degree (field, university, year) and research interests.
\end{IEEEbiography}

\begin{IEEEbiography}[{\includegraphics[width=1in,height=1.25in,clip,keepaspectratio]{photos/leu.jpg}}]{Jenq-Shiou Leu}
(Senior Member, IEEE) received the B.S. degree in mathematics and the M.S. degree in computer science and information engineering from National Taiwan University, Taipei, Taiwan, and the Ph.D. degree in computer science from National Tsing Hua University, Hsinchu, Taiwan.
From 1995 to 2007 he worked in industry, with Rising Star Technology, Mobitai Communications and Taiwan Mobile.
He joined the Department of Electronic and Computer Engineering, National Taiwan University of Science and Technology (NTUST), Taipei, Taiwan, in 2007, where he has been a Professor since 2014, chaired the department from 2017 to 2020, and has been Dean of the College of Electrical Engineering and Computer Science since 2024.
He leads the Mobilizing Information Technology Laboratory (MITLab).
His research interests include B5G/6G technology, AI/ML platforms, applications of large language models and biomedical information analysis.
\end{IEEEbiography}

\begin{IEEEbiography}[{\includegraphics[width=1in,height=1.25in,clip,keepaspectratio]{photos/ardiansyah.jpg}}]{Dhista Dwi Nur Ardiansyah}
received the bachelor's degree in informatics engineering education from Universitas Negeri Yogyakarta, Yogyakarta, Indonesia, in 2024.
Since 2026, Ardiansyah has been a Research Assistant with the Mobilizing Information Technology Laboratory (MITLab), National Taiwan University of Science and Technology, Taipei, Taiwan, through the Taiwan Experience Education Program (TEEP), working on small language models on edge devices, and has worked as a software engineer in Indonesia since 2021.
As an undergraduate, Ardiansyah won a gold medal at the 2021 Indonesian student digital innovation competition (LIDM) for an automatic question generator for teachers.
\end{IEEEbiography}
 in main.tex (after \bibliography).
% Author biographies (supervisor, 2026-11-24). Photos in manuscript/photos/ (git-ignored: the repo is public).
% Leu: degrees and career from his NTUST profile (et.ntust.edu.tw, fetched 2026-11-23) and the biography in his
% IEEE papers; photo from the ICOMPAC 2025 keynote page (icompac.its.ac.id), the largest found (1101x1357).
% Ardiansyah: from the owner's CV. Yunanto: from the owner (2026-11-24).
\begin{IEEEbiography}[{\includegraphics[width=1in,height=1.25in,clip,keepaspectratio]{photos/yunanto.jpg}}]{Rifky Ari Yunanto}
is pursuing the M.S. degree with the College of Electrical Engineering and Computer Science, National Taiwan University of Science and Technology (NTUST), Taipei, Taiwan, where he is a member of the Mobilizing Information Technology Laboratory (MITLab).
% TODO: Rifky's bachelor's degree (field, university, year) and research interests.
\end{IEEEbiography}

\begin{IEEEbiography}[{\includegraphics[width=1in,height=1.25in,clip,keepaspectratio]{photos/leu.jpg}}]{Jenq-Shiou Leu}
(Senior Member, IEEE) received the B.S. degree in mathematics and the M.S. degree in computer science and information engineering from National Taiwan University, Taipei, Taiwan, and the Ph.D. degree in computer science from National Tsing Hua University, Hsinchu, Taiwan.
From 1995 to 2007 he worked in industry, with Rising Star Technology, Mobitai Communications and Taiwan Mobile.
He joined the Department of Electronic and Computer Engineering, National Taiwan University of Science and Technology (NTUST), Taipei, Taiwan, in 2007, where he has been a Professor since 2014, chaired the department from 2017 to 2020, and has been Dean of the College of Electrical Engineering and Computer Science since 2024.
He leads the Mobilizing Information Technology Laboratory (MITLab).
His research interests include B5G/6G technology, AI/ML platforms, applications of large language models and biomedical information analysis.
\end{IEEEbiography}

\begin{IEEEbiography}[{\includegraphics[width=1in,height=1.25in,clip,keepaspectratio]{photos/ardiansyah.jpg}}]{Dhista Dwi Nur Ardiansyah}
received the bachelor's degree in informatics engineering education from Universitas Negeri Yogyakarta, Yogyakarta, Indonesia, in 2024.
Since 2026, Ardiansyah has been a Research Assistant with the Mobilizing Information Technology Laboratory (MITLab), National Taiwan University of Science and Technology, Taipei, Taiwan, through the Taiwan Experience Education Program (TEEP), working on small language models on edge devices, and has worked as a software engineer in Indonesia since 2021.
As an undergraduate, Ardiansyah won a gold medal at the 2021 Indonesian student digital innovation competition (LIDM) for an automatic question generator for teachers.
\end{IEEEbiography}
